% ============================================================
% content/solutions.tex —— 2024 湖南卷物理 · 参考答案与详细解析
% 解析类型：source_embedded（转录自来源站 Word 版内嵌解析，独立复核）
% ============================================================

\begin{answer}
\gksolhead{1}{单项选择题}
\gkanswer{答案 B}
普朗克认为黑体辐射的能量是量子化的，故 A 错误。入射光的频率大于或等于金属的截止频率才能产生光电效应；紫光的频率大于红光的频率，红光能让金属产生光电效应，则紫光也能让该金属产生光电效应，故 B 正确。石墨对 X 射线的散射过程遵循动量守恒，光子和电子碰撞后，电子获得一定的动量，光子动量变小，根据 $\lambda=\dfrac{h}{p}$ 可知波长变长，故 C 错误。德布罗意认为物质（包括电子）都具有波动性，故 D 错误。
\end{answer}

\begin{answer}
\gksolhead{2}{单项选择题}
\gkanswer{答案 D}
本题考查波长、频率和波速的关系及波的图像的理解。由图可知
\[
  x_{AB}=\frac{3}{2}\lambda=\SI{6}{\metre},
\]
得 $\lambda=\SI{4}{\metre}$，故 A 错误。波源的振动频率为
\[
  f=\frac{60}{60}\,\si{\hertz}=\SI{1}{\hertz},
\]
所以波速为 $v=\lambda f=\SI{4}{\metre\per\second}$，故 B 错误。振动周期为 $T=\SI{1}{\second}$，$0.25\,\si{\second}=\dfrac{T}{4}$，故 $B$ 点在 $t_{0}+0.25\,\si{\second}$ 时运动到平衡位置，位移为 0、速度最大，故 C 错误；$0.5\,\si{\second}=\dfrac{T}{2}$，故 $A$ 点在 $t_{0}+0.5\,\si{\second}$ 时运动到波峰，位移最大、速度为 0，故 D 正确。

\end{answer}

\begin{answer}
\gksolhead{3}{单项选择题}
\gkanswer{答案 A}
本题考查牛顿第二定律的瞬时性问题。剪断前，对 $B$、$C$、$D$ 整体分析得
\[
  F_{AB}=(3m+2m+m)g,
\]
对 $D$ 有 $F_{CD}=mg$。剪断后，对 $B$ 有
\[
  F_{AB}-3mg=3ma_{B},
\]
解得 $a_{B}=g$，方向竖直向上；对 $C$ 有
\[
  F_{CD}+2mg=2ma_{C},
\]
解得 $a_{C}=1.5g$，方向竖直向下。故选项 A 正确。
\end{answer}

\begin{answer}
\gksolhead{4}{单项选择题}
\gkanswer{答案 C}
本题考查导体切割磁感线产生感应电动势。如图所示，该导线在纸面内绕 $O$ 点逆时针转动，相当于 $Oa$、$Ob$、$Oc$ 三根导体棒转动切割磁感线，根据右手定则可知 $O$ 点电势最高。根据
\[
  E=Blv=\frac{1}{2}B\omega l^{2},
\]
又 $l_{Ob}=l_{Oc}=\sqrt{5}\,R>l_{Oa}$，所以 $0<U_{Oa}<U_{Ob}=U_{Oc}$，得
\[
  \varphi_{O}>\varphi_{a}>\varphi_{b}=\varphi_{c},
\]
故选项 C 正确。

\gkfig{figures/s4-rotate}
\end{answer}

\begin{answer}
\gksolhead{5}{单项选择题}
\gkanswer{答案 D}
解法一　根据异种点电荷的电场分布情况，可知两个点电荷的连线上，在两电荷中间及电荷量小的点电荷外侧，均有电势为零的点，所以 $x$ 正半轴上各点电势 $\varphi$ 应先负后正，由此可判断选项 D 正确。

解法二　根据点电荷周围的电势公式 $\varphi=k\dfrac{Q}{r}$（该公式教材中没有），设 $x'$ 处（$x'>0$）的电势为 0，得
\[
  k\frac{4q}{1+x'}+k\frac{-q}{x'}=0,
\]
解得 $x'=\dfrac{1}{3}$，故可知当 $0<x<\dfrac{1}{3}$ 时，$\varphi<0$；当 $x>\dfrac{1}{3}$ 时，$\varphi>0$。故选项 D 正确。
\end{answer}

\begin{answer}
\gksolhead{6}{单项选择题}
\gkanswer{答案 A}
本题考查交变电流产生和远距离输电原理，涉及等效法。转子在磁场中转动时产生的电动势的有效值为
\[
  E=\frac{NBS\omega}{\sqrt{2}},
\]
当转子角速度增加一倍时，升压变压器原副线圈两端电压都增加一倍，输电线上的电流变为 $I_{2}'=2I_{2}$，故 $R_{0}$ 上消耗的电功率变为原来的 4 倍，故 A 正确。升压变压器副线圈匝数增加一倍，副线圈两端电压增加一倍，输电线上的电流增加一倍，故 $R_{0}$ 上消耗的电功率变为原来的 4 倍，故 C 错误。

设降压变压器的原副线圈匝数比为 $k:1$，将降压变压器和用电器整体等效成一个电阻，等效电阻为
\[
  \frac{U_{3}}{I_{3}}=\frac{kU_{4}}{\frac{1}{k}I_{4}}=k^{2}R,
\]
等效电路图如图所示，则输电线上的电流为
\[
  I_{2}=\frac{U_{2}}{R_{0}+k^{2}R},
\]
$R_{0}$ 消耗的功率 $P=I_{2}^{2}R_{0}$。若 $R_{0}$ 阻值增加一倍，输电线上的电流
\[
  I_{2}''=\frac{U_{2}}{2R_{0}+k^{2}R},
\]
$R_{0}$ 消耗的功率 $P_{1}=I_{2}''^{2}\cdot 2R_{0}\ne 4P$，故 B 错误。若在用户端再并联一个完全相同的电阻 $R$，用户端电阻减为原来的一半，输电线上的电流
\[
  I_{2}'''=\frac{U_{2}}{R_{0}+k^{2}\frac{R}{2}}=\frac{2U_{2}}{2R_{0}+k^{2}R},
\]
$R_{0}$ 消耗的功率 $P_{2}=I_{2}'''^{2}R_{0}\ne 6P$，故 D 错误。

\gkfig{figures/s6-equiv}
\end{answer}

\begin{answer}
\gksolhead{7}{多项选择题}
\gkanswer{答案 BD}
本题考查万有引力定律、宇宙速度的理解。返回器绕月运行的轨道半径近似为月球半径，则有
\[
  G\frac{m_{\text{月}}m}{r_{\text{月}}^{2}}=m\frac{v_{\text{月}}^{2}}{r_{\text{月}}},
\]
在月球表面，由万有引力和重力的关系有
\[
  G\frac{m_{\text{月}}m'}{r_{\text{月}}^{2}}=m'g_{\text{月}},
\]
解得 $v_{\text{月}}=\sqrt{g_{\text{月}}r_{\text{月}}}$；同理可得地球第一宇宙速度即近地卫星的环绕速度 $v_{\text{地}}=\sqrt{g_{\text{地}}r_{\text{地}}}$，代入题中数据可得
\[
  v_{\text{月}}=\frac{\sqrt{6}}{12}v_{\text{地}},
\]
故 A 错误，B 正确。根据 $T=\dfrac{2\pi r}{v}$，得
\[
  T_{\text{月}}=\sqrt{\frac{3}{2}}\,T_{\text{地}},
\]
故 C 错误，D 正确。
\end{answer}

\begin{answer}
\gksolhead{8}{多项选择题}
\gkanswer{答案 CD}
本题考查电磁感应的能量、动量问题，涉及微元法。设平行金属导轨间距为 $L$，金属杆切割磁感线有 $E=BLv$，而 $I=\dfrac{E}{2R}$；根据动量定理，金属杆在 $AA_{1}B_{1}B$ 区域运动的过程中有
\[
  -BIL\,\Delta t=m\Delta v,
\]
即可得 $-BLq=mv_{B}-mv_{0}$，而
\[
  q=\frac{\Delta\varphi}{2R}=\frac{BLd}{2R},
\]
所以 $v_{B}=v_{0}-\dfrac{B^{2}L^{2}d}{2mR}$；设金属杆在 $BB_{1}C_{1}C$ 区域运动的时间为 $t_{0}$，同理可得，金属杆在 $BB_{1}C_{1}C$ 区域运动的过程中有
\[
  -\frac{B^{2}L^{2}d}{2R}-\mu mgt_{0}=-mv_{B},
\]
解得 $v_{B}=\dfrac{B^{2}L^{2}d}{2mR}+\mu gt_{0}$，联立可得
\[
  v_{B}=\frac{v_{0}}{2}+\frac{\mu gt_{0}}{2}>\frac{v_{0}}{2},
\]
则金属杆经过 $BB_{1}$ 的速度大于 $\dfrac{v_{0}}{2}$，故 A 错误。

在整个过程中，根据能量守恒定律有
\[
  \frac{1}{2}mv_{0}^{2}=\mu mgd+Q,
\]
则在整个过程中，定值电阻 $R$ 产生的热量为
\[
  Q_{R}=\frac{1}{2}Q=\frac{1}{4}mv_{0}^{2}-\frac{1}{2}\mu mgd,
\]
故 B 错误。

金属杆经过 $AA_{1}B_{1}B$ 与 $BB_{1}C_{1}C$ 区域，金属杆所受安培力的冲量为
\[
  -\sum BIL\,\Delta t=-\sum\frac{B^{2}L^{2}}{2R}\,v_{t}\Delta t=-\frac{B^{2}L^{2}x}{2R},
\]
金属杆经过 $AA_{1}B_{1}B$ 与 $BB_{1}C_{1}C$ 区域滑行距离均为 $d$，则金属杆所受安培力的冲量相同，故 C 正确。

根据 A 选项可知，金属杆以初速度 $v_{0}$ 在磁场中运动有
\[
  -\frac{B^{2}L^{2}\times 2d}{2R}-\mu mgt_{0}=-mv_{0};
\]
设金属杆的初速度加倍后，金属杆通过 $BB_{1}C_{1}C$ 区域的时间为 $t_{1}$，金属杆在磁场中运动的距离为 $x$，则有
\[
  -\frac{B^{2}L^{2}x}{2R}-\mu mgt_{1}=0-2mv_{0},
\]
联立解得
\[
  x=\frac{2R}{B^{2}L^{2}}(2mv_{0}-\mu mgt_{1})
    =\frac{2mv_{0}-\mu mgt_{1}}{mv_{0}-\mu mgt_{0}}\times 2d,
\]
由于 $t_{1}<t_{0}$，所以 $\dfrac{2mv_{0}-\mu mgt_{1}}{mv_{0}-\mu mgt_{0}}>2$，则 $x>4d$，即金属杆在磁场中运动的距离大于原来的 2 倍，故 D 正确。
\end{answer}

\begin{answer}
\gksolhead{9}{多项选择题}
\gkanswer{答案 BC}
本题考查双缝干涉的理解。当 $\theta=0^{\circ}$ 时，画出光路图，如图甲所示，$S$ 发出的光与通过平面镜反射的光（可以等效成虚像 $S'$ 发出的光）是同一列光分成的，满足相干光条件，且该干涉可看成双缝干涉，设 $S$ 与 $S'$ 的距离为 $d$，$S$ 到光屏的距离为 $l$，代入双缝干涉条纹间距公式可得
\[
  \Delta x=\frac{l\lambda}{d},
\]
则沿 $OA$ 向右略微平移平面镜，对 $l$ 和 $d$ 均没有影响，干涉条纹间距不变，故 C 正确；若沿 $AO$ 向左略微平移平面镜，由图甲可以看出，干涉条纹向远离 $A$ 处移动，故 D 错误。

\gkfig{figures/s9-case0}

当 $\theta\ne 0^{\circ}$ 时，作出光路图如图乙所示，沿 $OA$ 向右略微平移平面镜，即图乙中从①位置到②位置，由图乙可看出双缝的间距增大，则干涉条纹间距减小，故 B 正确；沿 $AO$ 向左略微平移平面镜，即图乙中从②位置到①位置，由图乙可看出干涉条纹向上移动，故 A 错误。

\gkfig{figures/s9-tilt}
\end{answer}

\begin{answer}
\gksolhead{10}{多项选择题}
\gkanswer{答案 AC}
两小球均做匀速直线运动，因为 $A$、$B$ 两小球能相遇，可绘出运动轨迹图如图所示，根据正弦定理有
\[
  \frac{v_{2}t_{1}}{\sin(2\alpha+\theta)}
  =\frac{v_{2}'t_{2}}{\sin\theta}
  =\frac{(t_{1}+t_{2})v_{1}}{\sin[180^{\circ}-(2\alpha+2\theta)]},
\]
若 $\theta=15^{\circ}$，代入数据 $v_{2}=2$~m/s，$v_{2}'=1$~m/s 解得，当 $\alpha=15^{\circ}$ 时，$v_{1}$ 取得最大值为 $\sqrt{2}$~m/s，故 A 正确，B 错误；若 $\theta=30^{\circ}$，代入数据 $v_{2}=2$~m/s，$v_{2}'=1$~m/s 解得，当 $\alpha=0^{\circ}$ 时，$v_{1}$ 取得最大值为 $\dfrac{2}{3}\sqrt{3}$~m/s，故 C 正确，D 错误。

\gkfig{figures/s10-trajectory}
\end{answer}

\begin{answer}
\gksolhead{11}{实验题}
\gkanswer{答案 (1) 1.23 V，12.3 Ω；(4) 右；(5) 小于}
本题考查电学创新实验。(1) 电压表分度值为 \SI{0.1}{\volt}，则电压表读数需估读一位，读数为 \SI{1.23}{\volt}，根据欧姆定律可知，金属丝的测量值
\[
  R=\frac{U}{I}=\SI{12.3}{\ohm}.
\]

(4) 根据图丁可知气压越小电阻越大，则压强 $p$ 减小，则电阻 $R$ 增大，再根据 $U=IR$，故电压增大，电压表的指针位置应该在题图丙中指针位置的右侧。

(5) 电流表采用外接法会导致电压表分流，即 $R_{\text{测}}=\dfrac{U}{I_{\text{测}}}$，$I_{\text{测}}=I_{\text{真}}+\dfrac{U}{R_{V}}$，即 $I_{\text{测}}$ 偏大，故 $R_{\text{测}}<R_{\text{真}}$。
\end{answer}

\begin{answer}
\gksolhead{12}{实验题}
\gkanswer{答案 (4) 见图线；(5) 线性的；(6) 0.120 kg；(7) 负方向}
本题利用弹簧振子考查力学创新实验。

(4) 根据表格中数据，在图中描点并作出 $T^{2}$--$m$ 关系图线，如图所示。

\gkfig{figures/s12-tm}

(5) 图线是一条倾斜的直线，说明弹簧振子振动周期的二次方与砝码质量为线性关系。

(6) 在图线上找到 $T^{2}=\SI{0.880}{\second\squared}$ 的点，对应横坐标为 \SI{0.120}{\kilogram}。

(7) 换一个质量较小的滑块做实验，滑块和砝码总质量较原来的偏小，要得到相同的周期，应放质量更大的砝码，对应纵坐标点应右移，则所得图线与原图线相比下移，即沿纵轴负方向移动。
\end{answer}

\begin{answer}
\gksolhead{13}{解答题}
\gkanswer{答案 (1) $V_{0}=\dfrac{pV}{p_{0}}$；(2) $V=5\times10^{-3}$ m$^{3}$}
本题考查气体实验定律。

(1) 根据玻意耳定律有
\[
  pV=p_{0}V_{0},
\]
解得
\[
  V_{0}=\frac{pV}{p_{0}}.
\]

(2) 设气球内气体质量为 $m_{\text{气}}$，则 $m_{\text{气}}=\rho V$。根据阿基米德原理，气球受到的浮力为 $\rho_{0}gV$。对气球进行受力分析如图所示。

\gkfig{figures/s13-forces}

根据气球的受力分析有
\[
  mg+\rho_{0}gV=m_{\text{气}}g+m_{0}g.
\]
第 (1) 问中，气体等温变化时，气体质量不变，有
\[
  m_{\text{气}}=\rho V=\rho_{0}V_{0}.
\]
又因为 $pV=p_{0}V_{0}$，可得
\[
  \frac{p}{\rho}=\frac{p_{0}}{\rho_{0}}.
\]
结合题中 $p$ 和 $V$ 满足的关系为
\[
  (p-p_{0})(V-V_{B0})=C,
\]
解得 $V=5\times10^{-3}$~m$^{3}$。
\end{answer}

\begin{answer}
\gksolhead{14}{解答题}
\gkanswer{答案 (1) $B_{\min}=\dfrac{2\pi mv_{0}}{eL}$；(2) $|\tan\theta|=\dfrac{2\pi r}{L}$；(3) $y_{m}=\dfrac{2\pi^{2}r^{2}v_{0}^{2}m}{EeL^{2}}$}
(1) 电子在匀强磁场中运动时，将其分解为沿 $x$ 轴的匀速直线运动和在 $yOz$ 平面内的匀速圆周运动，设电子入射时沿 $y$ 轴的分速度大小为 $v_{y}$，由电子在 $x$ 轴方向做匀速直线运动得
\[
  L=v_{0}t.
\]
在 $yOz$ 平面内，设电子做匀速圆周运动的半径为 $R$，周期为 $T$，由牛顿第二定律得
\[
  Bev_{y}=m\frac{v_{y}^{2}}{R},
\]
可得
\[
  R=\frac{mv_{y}}{Be},\qquad T=\frac{2\pi R}{v_{y}}=\frac{2\pi m}{Be}.
\]
由题意可知所有电子均能经过 $O$ 进入电场，则有
\[
  t=nT\qquad(n=1,2,3,\dots)
\]
联立得
\[
  B=\frac{2\pi nmv_{0}}{eL},
\]
当 $n=1$ 时，$B$ 有最小值，可得
\[
  B_{\min}=\frac{2\pi mv_{0}}{eL}.
\]

(2) 将电子的速度分解，如图所示。

\gkfig{figures/s14-decomp}

有
\[
  \tan\theta=\frac{v_{y}}{v_{0}}.
\]
当 $\tan\theta$ 有最大值时，$v_{y}$ 最大，$R$ 最大，此时 $R=r$。又
\[
  B=\frac{2\pi mv_{0}}{eL},\qquad R=\frac{mv_{y}}{Be},
\]
联立可得
\[
  v_{ym}=\frac{2\pi v_{0}r}{L},\qquad \tan\theta=\frac{2\pi r}{L}.
\]

(3) 当 $v_{y}$ 最大时，电子在电场中运动时沿 $y$ 轴正方向有最大位移 $y_{m}$，根据匀变速直线运动规律有
\[
  y_{m}=\frac{v_{ym}^{2}}{2a},
\]
由牛顿第二定律知
\[
  a=\frac{Ee}{m},
\]
又 $v_{ym}=\dfrac{2\pi v_{0}r}{L}$，联立得
\[
  y_{m}=\frac{2\pi^{2}r^{2}v_{0}^{2}m}{EeL^{2}}.
\]
\end{answer}

\begin{answer}
\gksolhead{15}{解答题}
\gkanswer{答案 (1) $v=\dfrac{m_{A}v_{0}}{m_{A}+m_{B}}$，$F=\dfrac{m_{A}^{2}v_{0}^{2}}{(m_{A}+m_{B})R}$；(2) 2 或 5}
\gkanswer{(3) $s=\dfrac{2\pi Rm_{A}}{m_{A}+m_{B}}\cdot\dfrac{e^{2n}-1}{e^{2n}(e-1)}$}
本题考查多次碰撞的综合问题。

(1) 设碰撞后两小球的速度大小为 $v$，根据动量守恒定律有
\[
  m_{A}v_{0}=(m_{A}+m_{B})v,
\]
可得
\[
  v=\frac{m_{A}v_{0}}{m_{A}+m_{B}}.
\]
碰撞后根据牛顿第二定律有
\[
  F=(m_{A}+m_{B})\frac{v^{2}}{R},
\]
可得
\[
  F=\frac{m_{A}^{2}v_{0}^{2}}{(m_{A}+m_{B})R}.
\]

(2) 若两小球发生弹性碰撞，设碰后速度分别为 $v_{A}$、$v_{B}$，则碰后动量守恒和能量守恒有
\[
  m_{A}v_{0}=m_{A}v_{A}+m_{B}v_{B},
\]
\[
  \frac{1}{2}m_{A}v_{0}^{2}=\frac{1}{2}m_{A}v_{A}^{2}+\frac{1}{2}m_{B}v_{B}^{2},
\]
联立解得
\[
  v_{A}=\frac{m_{A}-m_{B}}{m_{A}+m_{B}}v_{0},\qquad
  v_{B}=\frac{2m_{A}v_{0}}{m_{A}+m_{B}}.
\]
因为所有的碰撞位置刚好位于等边三角形的三个顶点，如图所示。

\gkfig{figures/s15-triangle}

① 若第二次碰撞发生在图中的 $b$ 点，则从第一次碰撞到第二次碰撞之间，$A$、$B$ 通过的路程之比为
\[
  \frac{1+3k_{1}}{4+3k_{1}}\qquad(k_{1}=0,1,2,3,\dots)
\]
则有
\[
  \frac{v_{A1}}{v_{B1}}=\frac{x_{A1}}{x_{B1}}=\frac{1+3k_{1}}{4+3k_{1}}\qquad(k_{1}=0,1,2,3,\dots)
\]
联立解得
\[
  \frac{m_{A1}}{m_{B1}}=\frac{4+3k_{1}}{2-3k_{1}}\qquad(k_{1}=0,1,2,3,\dots)
\]
由于两质量均为正数，故 $k_{1}=0$，即 $\dfrac{m_{A1}}{m_{B1}}=2$。对第二次碰撞，设 $A$、$B$ 碰撞后的速度大小分别为 $v_{A1}'$、$v_{B1}'$，则同样有
\[
  m_{A1}v_{A1}+m_{B1}v_{B1}=m_{A1}v_{A1}'+m_{B1}v_{B1}',
\]
\[
  \frac{1}{2}m_{A1}v_{A1}^{2}+\frac{1}{2}m_{B1}v_{B1}^{2}
  =\frac{1}{2}m_{A1}v_{A1}'^{2}+\frac{1}{2}m_{B1}v_{B1}'^{2},
\]
联立解得 $v_{A1}'=v_{0}$，$v_{B1}'=0$，故第三次碰撞发生在 $b$ 点、第四次碰撞发生在 $c$ 点，以此类推，满足题意。

② 若第二次碰撞发生在图中的 $c$ 点，则从第一次碰撞到第二次碰撞之间，$A$、$B$ 通过的路程之比为
\[
  \frac{2+3k_{2}}{5+3k_{2}}\qquad(k_{2}=0,1,2,3,\dots)
\]
则有
\[
  \frac{v_{A2}}{v_{B2}}=\frac{x_{A2}}{x_{B2}}=\frac{2+3k_{2}}{5+3k_{2}}\qquad(k_{2}=0,1,2,3,\dots)
\]
联立可得
\[
  \frac{m_{A2}}{m_{B2}}=\frac{5+3k_{2}}{1-3k_{2}}\qquad(k_{2}=0,1,2,3,\dots)
\]
因为两质量均为正数，故 $k_{2}=0$，即 $\dfrac{m_{A2}}{m_{B2}}=5$。根据①的分析可证 $v_{A2}'=v_{0}$，$v_{B2}'=0$，故第三次碰撞发生在 $b$ 点、第四次碰撞发生在 $c$ 点，以此类推，满足题意。

综上可知 $\dfrac{m_{A}}{m_{B}}=2$ 或 5。

(3) 第一次碰前相对速度大小为 $v_{0}$，第一次碰后的相对速度大小为
\[
  v_{1\text{相}}=ev_{0},
\]
第一次碰后与第二次相碰前 $B$ 比 $A$ 多运动一圈，即 $B$ 相对 $A$ 运动一圈，有
\[
  t_{1}=\frac{2\pi R}{v_{1\text{相}}}.
\]
第一次碰撞动量守恒有
\[
  m_{A}v_{0}=m_{A}v_{A}'+m_{B}v_{B}',
\]
且
\[
  v_{1\text{相}}=v_{B}'-v_{A}'=ev_{0},
\]
联立解得
\[
  v_{B}'=\frac{m_{A}}{m_{A}+m_{B}}(v_{0}+v_{1\text{相}}).
\]
$B$ 运动的路程
\[
  s_{1}=v_{B}'t_{1}
     =\frac{2\pi Rm_{A}}{m_{A}+m_{B}}\left(\frac{v_{0}}{v_{1\text{相}}}+1\right)
     =\frac{2\pi Rm_{A}}{m_{A}+m_{B}}\left(\frac{1}{e}+1\right).
\]
第二次碰撞的相对速度大小为
\[
  v_{2\text{相}}=ev_{1\text{相}}=e^{2}v_{0},
\]
\[
  t_{2}=\frac{2\pi R}{v_{2\text{相}}}.
\]
第二次碰撞有
\[
  m_{A}v_{0}=m_{A}v_{A}''+m_{B}v_{B}'',
\]
且
\[
  v_{2\text{相}}=v_{A}''-v_{B}''=e^{2}v_{0},
\]
联立可得
\[
  v_{B}''=\frac{m_{A}}{m_{A}+m_{B}}(v_{0}-v_{2\text{相}}).
\]
所以 $B$ 运动的路程
\[
  s_{2}=v_{B}''t_{2}
     =\frac{2\pi Rm_{A}}{m_{A}+m_{B}}\left(\frac{v_{0}}{v_{2\text{相}}}-1\right)
     =\frac{2\pi Rm_{A}}{m_{A}+m_{B}}\left(\frac{1}{e^{2}}-1\right).
\]
一共碰了 $2n$ 次，有
\[
  s=s_{1}+s_{2}+s_{3}+\cdots+s_{2n}
   =\frac{2\pi Rm_{A}}{m_{A}+m_{B}}\left(\frac{1}{e}+\frac{1}{e^{2}}+\frac{1}{e^{3}}+\cdots+\frac{1}{e^{2n}}\right)
   =\frac{2\pi Rm_{A}}{m_{A}+m_{B}}\cdot\frac{e^{2n}-1}{e^{2n}(e-1)}.
\]
\end{answer}
